% TOP_PROBLEM: {"id": 1100208, "title": "Questions in Geometric Group Theory — Q 2.8", "category_id": 17, "category": {"id": 17, "name": "group_theory", "display_name": "Group Theory", "description": "Problems about groups, group actions, representations, and related algebraic structures.", "slug": "group-theory", "order_index": 17, "created_at": "2026-07-31T15:26:25.671Z"}, "status": "open", "problem_number": "AMR-010-0208", "set_id": 13, "statusQualification": "Uses the all-subgroups convention for the Tits alternative, explicitly given by Bestvina--Feighn--Handel, rather than restricting the question to finitely generated subgroups."}
% FIRST_POSTED: 2026-10-01
% ENTRY_KIND: statement_only
% UnsolvedMath ID 1100208; source catalogue code AMR-010-0208
% !TeX program = lualatex
% Definition and sources only; this file is not an LLM solution attempt.
\documentclass[11pt,a4paper]{article}
\usepackage[margin=25mm]{geometry}
\usepackage{fontspec}
\setmainfont{lmroman10-regular.otf}[BoldFont=lmroman10-bold.otf,
 ItalicFont=lmroman10-italic.otf,BoldItalicFont=lmroman10-bolditalic.otf]
\usepackage{amsmath,amssymb,mathtools}
\usepackage{unicode-math}
\setmathfont{latinmodern-math.otf}
\AtBeginDocument{\let\setminus\smallsetminus}
\usepackage{xurl,hyperref}
\hypersetup{colorlinks=true,urlcolor=blue}
\setcounter{secnumdepth}{0}
\setlength{\parindent}{0pt}
\setlength{\parskip}{5pt}
\setlength{\emergencystretch}{3em}
\begin{document}
\section*{Questions in Geometric Group Theory — Q 2.8}\label{1100208}
{\small UnsolvedMath ID 1100208 · AMR-010-0208 · Group Theory · AMR Open Problem Lists}
\subsection{Definitions and mathematical statement}
A group satisfies the Tits alternative if every subgroup either contains a subgroup isomorphic to the free group $F_2$ or has a solvable subgroup of finite index. Solvable means that repeatedly taking commutator subgroups eventually gives the trivial group. The finite-index solvable subgroup need not be normal, although a normal finite-index subgroup can be obtained by taking its core.

A CAT$(0)$ group admits a proper cocompact isometric action on a proper geodesic space satisfying Euclidean triangle comparison. An automatic structure for a finitely generated group consists of a regular language $L$ of words representing every element, such that for every generator $s$, and also for equality, the padded pairs
\[
 \{(u,v)\in L^2:\overline u s=\overline v\}
\]
are recognized by finite synchronous automata. A regular language is one recognized by a finite-state automaton; padding aligns the ends of unequal-length words. Biautomaticity additionally requires the analogous automata for $s\overline u=\overline v$.

Do all CAT$(0)$ groups satisfy the Tits alternative? Ask the same question separately for automatic and for biautomatic groups. The subgroup being tested need not inherit the ambient automatic structure or act cocompactly on the same CAT$(0)$ space. The subgroup need not be finitely generated. This is the all-subgroups convention explicitly used by Bestvina--Feighn--Handel [S3]. Testing only finitely generated subgroups would impose a potentially weaker requirement. The three ambient classes are separate versions of the source question.
\subsection{Short English statement}
Do CAT(0), automatic, and biautomatic groups satisfy the Tits alternative for every subgroup?
\subsection{Sources}
\begin{itemize}
\item[S1] ULAM.AI and the UnsolvedMath Contributors, supplied problems.json, record 1100208 (AMR-010-0208), AMR Open Problem Lists. Catalogue identity and source transcription; CC BY 4.0.
\url{https://huggingface.co/datasets/ulamai/UnsolvedMath}.
\item[S2] Bestvina - Questions in Geometric Group Theory (pdf) (2004), Question 2.8 (PDF page 7). Original source indexed by AMR-010-0208.
\url{https://www.math.utah.edu/~bestvina/eprints/questions-updated.pdf}.
\item[S3] M. Bestvina, M. Feighn and M. Handel, The Tits alternative for Out($F_n$) I, Annals of Mathematics 151 (2000), p. 518. Definition for all subgroups.
\url{https://emis.de/ft/50655}.
\end{itemize}
\end{document}
